% Pandoc LaTeX template — hỗ trợ tiếng Việt (XeLaTeX)
\documentclass[12pt,a4paper]{article}

% Font & ngôn ngữ
\usepackage{fontspec}
\usepackage{polyglossia}
\setmainlanguage{vietnamese}
\setotherlanguage{english}
\defaultfontfeatures{Ligatures=TeX}
% Latin Modern được cài cùng TeX Live, có đủ glyph tiếng Việt và không phụ
% thuộc font hệ điều hành. Code dùng font mono thật để sơ đồ/căn cột không lệch.
\setmainfont{lmsans10-regular.otf}[
  BoldFont = lmsans10-bold.otf,
  ItalicFont = lmsans10-oblique.otf,
  BoldItalicFont = lmsans10-boldoblique.otf
]
\setsansfont{lmsans10-regular.otf}[
  BoldFont = lmsans10-bold.otf,
  ItalicFont = lmsans10-oblique.otf,
  BoldItalicFont = lmsans10-boldoblique.otf
]
\setmonofont{lmmonolt10-regular.otf}[
  BoldFont = lmmonolt10-bold.otf,
  ItalicFont = lmmonolt10-oblique.otf,
  BoldItalicFont = lmmonolt10-boldoblique.otf,
  Scale = MatchLowercase
]

% Layout
\usepackage[margin=2.5cm, top=2.8cm, bottom=2.8cm]{geometry}
\usepackage{setspace}
\onehalfspacing
\setlength{\emergencystretch}{3em}
\raggedbottom
\usepackage{microtype}

% Màu & link
\usepackage{xcolor}
\definecolor{linkcolor}{RGB}{45,90,30}
\definecolor{codebg}{RGB}{245,245,245}
\definecolor{accent}{RGB}{80,130,40}
\usepackage{hyperref}
\hypersetup{
  colorlinks=true,
  linkcolor=linkcolor,
  urlcolor=linkcolor,
  citecolor=linkcolor,
  pdfauthor={FitLife - Module Người 4},
  pdftitle={Đặc tả Dashboard, Notification và Badges}
}

% Bảng
\usepackage{longtable,booktabs,array,calc}
\usepackage{tabularx}
\usepackage{etoolbox}
\renewcommand{\arraystretch}{1.3}
\setlength{\tabcolsep}{4pt}
\AtBeginEnvironment{longtable}{\small}
% Fix pandoc 3.x table caption counter
\newcounter{none}
\renewcommand{\thenone}{}

% Code listing — idiomatic mode dùng lstlisting (tự xuống dòng, không tràn lề)
\usepackage{listings}
\lstset{
  basicstyle=\ttfamily\footnotesize,
  backgroundcolor=\color{codebg},
  breaklines=true,
  breakatwhitespace=false,
  frame=single,
  rulecolor=\color{gray!40},
  columns=fullflexible,
  keepspaces=true,
  showstringspaces=false,
  tabsize=2,
  xleftmargin=0.5em,
  xrightmargin=0.5em,
  aboveskip=1em,
  belowskip=1em,
}
% Pandoc 3.10 bọc inline code bằng \passthrough khi dùng idiomatic listings.
\providecommand{\passthrough}[1]{#1}

% Mục lục — dùng mặc định của LaTeX (không cần package thêm)

% Pandoc list helper
\providecommand{\tightlist}{%
  \setlength{\itemsep}{0pt}\setlength{\parskip}{0pt}}

% Pandoc variables
\title{Tài liệu Giải thích Module Dashboard, Notification và Badges}
\author{FitLife --- Người 4}
\date{\today}

\begin{document}
\maketitle
\tableofcontents
\newpage

\textbf{Dự án:} FitLife --- Ứng dụng tập luyện \& sức khỏe\\
\textbf{Phạm vi:} Module của Người 4 (theo bản phân công nhóm)\\
\textbf{Công nghệ:} React Native (Expo) + Supabase (PostgreSQL)

\begin{center}\rule{0.5\linewidth}{0.5pt}\end{center}

\subsection{Module này làm gì?}\label{module-nuxe0y-luxe0m-guxec}

Module \textbf{Dashboard, Notification \& Badges} là phần giúp người
dùng \textbf{nhìn tổng quan tiến độ}, \textbf{nhận nhắc nhở} và
\textbf{được ghi nhận thành tích} khi tập luyện đều đặn.

Gồm \textbf{3 chức năng chính:}

{\def\LTcaptype{none} % do not increment counter
\begin{longtable}[]{@{}
  >{\raggedright\arraybackslash}p{(\linewidth - 4\tabcolsep) * \real{0.2750}}
  >{\raggedright\arraybackslash}p{(\linewidth - 4\tabcolsep) * \real{0.4750}}
  >{\raggedright\arraybackslash}p{(\linewidth - 4\tabcolsep) * \real{0.2500}}@{}}
\toprule\noalign{}
\begin{minipage}[b]{\linewidth}\raggedright
Chức năng
\end{minipage} & \begin{minipage}[b]{\linewidth}\raggedright
Người dùng thấy gì
\end{minipage} & \begin{minipage}[b]{\linewidth}\raggedright
Mục đích
\end{minipage} \\
\midrule\noalign{}
\endhead
\bottomrule\noalign{}
\endlastfoot
\textbf{Dashboard} & Tab Home: chuỗi tập, calo, số buổi tập, biểu đồ 7
ngày, sức khỏe và huy hiệu & Tổng hợp số liệu quan trọng trên một màn
hình \\
\textbf{Notification} & Thông báo nhắc uống nước, nhắc tập, thông báo
huy hiệu mới & Giúp user không quên tập và uống nước \\
\textbf{Badges} & Lưới huy hiệu (Khởi đầu, Kiên trì 3 ngày, Đốt cháy 500
kcal\ldots) & Gamification --- khuyến khích duy trì thói quen \\
\end{longtable}
}

\textbf{Người 4 phụ trách cả 3 lớp:} - \textbf{Frontend} --- giao diện
Dashboard, huy hiệu, cài đặt nhắc nhở - \textbf{Backend} --- API tổng
hợp Dashboard, bảng huy hiệu, lưu cài đặt thông báo - \textbf{Test} ---
kiểm thử thông báo, quyền OS, logic huy hiệu

\begin{center}\rule{0.5\linewidth}{0.5pt}\end{center}

\subsection{Kiến trúc tổng quan (dễ
hiểu)}\label{kiux1ebfn-truxfac-tux1ed5ng-quan-dux1ec5-hiux1ec3u}

Hãy tưởng tượng app như một nhà 3 tầng:

\begin{lstlisting}
+------------------------------------------+
|  TANG 1 - GIAO DIEN (man hinh user nhin) |
|  DashboardScreen, BadgeGrid, Modal nhac  |
+----------------------+-------------------+
                       | goi ham
+----------------------v-------------------+
|  TANG 2 - SERVICE (xu ly nghiep vu)      |
|  dashboardService, badgeService,         |
|  notificationService, waterReminder      |
+----------------------+-------------------+
                       | goi API / RPC
+----------------------v-------------------+
|  TANG 3 - SUPABASE (luu tru & bao mat)   |
|  PostgreSQL + RLS + RPC functions        |
+------------------------------------------+
\end{lstlisting}

\textbf{Quy tắc quan trọng:} - Màn hình \textbf{không} gọi database trực
tiếp --- luôn qua \textbf{service}. - Logic tính toán thuần (vd: đủ điều
kiện huy hiệu chưa?) nằm trong \textbf{\passthrough{\lstinline!lib/!}}
để dễ test. - Mỗi user \textbf{chỉ thấy dữ liệu của mình} nhờ RLS (Row
Level Security) trên Supabase.

\begin{center}\rule{0.5\linewidth}{0.5pt}\end{center}

\subsection{Dashboard (Trang chủ tổng
quan)}\label{dashboard-trang-chux1ee7-tux1ed5ng-quan}

\subsubsection{User thấy gì trên
Dashboard?}\label{user-thux1ea5y-guxec-truxean-dashboard}

Khi mở tab \textbf{Home}, màn hình hiển thị:

\begin{enumerate}
\def\labelenumi{\arabic{enumi}.}
\tightlist
\item
  \textbf{Lưới chỉ số nhanh} - streak hiện tại, calo đốt hôm nay, số
  buổi tập hôm nay và streak dài nhất.
\item
  \textbf{Biểu đồ 7 ngày} --- Hoạt động tập luyện trong tuần qua (mỗi
  cột = 1 ngày)
\item
  \textbf{Thẻ buổi tập} - hiện đang là dữ liệu trình diễn cố định
  (\passthrough{\lstinline!Power Strength II!}, 12/24, 50\%), chưa đọc
  chương trình hiện tại từ database.
\item
  \textbf{Sức khỏe hôm nay} - giờ nhắc tập, calo nạp, nước và macro dinh
  dưỡng.
\item
  \textbf{Huy hiệu} - các badge đã đạt (sáng) và chưa đạt (mờ).
\end{enumerate}

Kéo xuống để \textbf{làm mới} (pull-to-refresh). Sau khi hoàn thành buổi
tập, Dashboard \textbf{tự cập nhật} số liệu mới.

\subsubsection{Dữ liệu lấy từ
đâu?}\label{dux1eef-liux1ec7u-lux1ea5y-tux1eeb-ux111uxe2u}

Riêng phần \textbf{số liệu tổng quan}, app ưu tiên một RPC trên server
thay vì đọc từng bảng:

\textbf{\passthrough{\lstinline!get\_dashboard\_summary()!}} --- hàm
PostgreSQL trả về JSON gồm:

{\def\LTcaptype{none} % do not increment counter
\begin{longtable}[]{@{}
  >{\raggedright\arraybackslash}p{(\linewidth - 2\tabcolsep) * \real{0.4706}}
  >{\raggedright\arraybackslash}p{(\linewidth - 2\tabcolsep) * \real{0.5294}}@{}}
\toprule\noalign{}
\begin{minipage}[b]{\linewidth}\raggedright
Trường
\end{minipage} & \begin{minipage}[b]{\linewidth}\raggedright
Ý nghĩa
\end{minipage} \\
\midrule\noalign{}
\endhead
\bottomrule\noalign{}
\endlastfoot
\passthrough{\lstinline!display\_name!} & Tên hiển thị \\
\passthrough{\lstinline!current\_streak!} & Chuỗi ngày tập liên tiếp
hiện tại \\
\passthrough{\lstinline!longest\_streak!} & Kỷ lục chuỗi dài nhất \\
\passthrough{\lstinline!calories\_burned\_today!} & Calo đốt hôm nay \\
\passthrough{\lstinline!workouts\_today!} & Số buổi tập hôm nay \\
\passthrough{\lstinline!weekly\_workouts!} & Mảng 7 ngày: mỗi ngày có
\passthrough{\lstinline!date!}, \passthrough{\lstinline!workouts!},
\passthrough{\lstinline!calories\_burned!} \\
\passthrough{\lstinline!wakeup\_time!} & Giờ dậy (dùng cho nhắc tập) \\
Các cờ \passthrough{\lstinline!*\_reminder\_enabled!} & Đã bật nhắc
nước/tập/badge chưa \\
\end{longtable}
}

\textbf{Ví dụ JSON trả về:}

\begin{lstlisting}
{
  "display_name": "Minh Sang",
  "current_streak": 3,
  "longest_streak": 5,
  "calories_burned_today": 320,
  "workouts_today": 1,
  "weekly_workouts": [
    { "date": "2026-06-24", "workouts": 0, "calories_burned": 0 },
    { "date": "2026-06-25", "workouts": 1, "calories_burned": 280 },
    { "date": "2026-06-26", "workouts": 0, "calories_burned": 0 },
    { "date": "2026-06-27", "workouts": 1, "calories_burned": 310 },
    { "date": "2026-06-28", "workouts": 1, "calories_burned": 290 },
    { "date": "2026-06-29", "workouts": 0, "calories_burned": 0 },
    { "date": "2026-06-30", "workouts": 1, "calories_burned": 320 }
  ],
  "wakeup_time": "06:30",
  "water_reminder_enabled": true,
  "workout_reminder_enabled": false,
  "badge_notifications_enabled": true
}

\end{lstlisting}

Dashboard vẫn gọi thêm luồng huy hiệu và cài đặt thông báo song song; vì
vậy không nên hiểu là toàn bộ màn hình chỉ có đúng một request.

\subsubsection{Biểu đồ 7 ngày hoạt động thế
nào?}\label{biux1ec3u-ux111ux1ed3-7-nguxe0y-houx1ea1t-ux111ux1ed9ng-thux1ebf-nuxe0o}

\textbf{Vấn đề đã gặp:} Khi user chưa tập buổi nào, biểu đồ hiển thị cột
rất mờ --- trông như lỗi.

\textbf{Cách xử lý:} - Server luôn trả \textbf{đủ 7 ngày} (kể cả ngày
không tập = 0 buổi, 0 calo). - Client chuẩn hóa lại qua hàm
\passthrough{\lstinline!normalizeWeeklyWorkouts()!} --- khớp đúng ngày,
điền 0 nếu thiếu. - Khi chưa có buổi tập nào -\textgreater{} hiển thị
thông báo \textbf{``Chưa có buổi tập nào''} thay vì cột trống mờ. - Ngày
hôm nay được \textbf{highlight} màu xanh lá; nhãn CN-T7 luôn hiển thị.

\textbf{Lưu ý múi giờ:} RPC dùng \passthrough{\lstinline!current\_date!}
theo timezone của database, còn client tạo khung 7 ngày theo timezone
thiết bị. Nếu hai timezone khác nhau, dữ liệu quanh nửa đêm có thể lệch
một ngày. Cần thống nhất timezone hoặc truyền ngày local vào RPC nếu
muốn xử lý triệt để.

\subsubsection{Nếu RPC lỗi thì
sao?}\label{nux1ebfu-rpc-lux1ed7i-thuxec-sao}

App thử \textbf{phương án dự phòng (fallback):}
\passthrough{\lstinline!dashboardService!} đọc
\passthrough{\lstinline!profiles!},
\passthrough{\lstinline!user\_streaks!} và
\passthrough{\lstinline!workout\_sessions!} rồi gom số liệu giống RPC.
Fallback hữu ích khi RPC chưa tồn tại hoặc trả lỗi, nhưng không bảo đảm
cứu được trường hợp mất mạng vì các truy vấn dự phòng cũng cần Supabase.

\subsubsection{File code liên quan}\label{file-code-liuxean-quan}

{\def\LTcaptype{none} % do not increment counter
\begin{longtable}[]{@{}
  >{\raggedright\arraybackslash}p{(\linewidth - 2\tabcolsep) * \real{0.4000}}
  >{\raggedright\arraybackslash}p{(\linewidth - 2\tabcolsep) * \real{0.6000}}@{}}
\toprule\noalign{}
\begin{minipage}[b]{\linewidth}\raggedright
File
\end{minipage} & \begin{minipage}[b]{\linewidth}\raggedright
Vai trò
\end{minipage} \\
\midrule\noalign{}
\endhead
\bottomrule\noalign{}
\endlastfoot
\passthrough{\lstinline!src/screens/DashboardScreen.tsx!} & Màn hình
chính, gọi load dữ liệu \\
\passthrough{\lstinline!src/services/dashboardService.ts!} & Gọi RPC +
chuẩn hóa 7 ngày \\
\passthrough{\lstinline!src/components/StatsGrid.tsx!} & Thẻ sức khỏe:
giờ nhắc, dinh dưỡng, nước và macro \\
\passthrough{\lstinline!src/components/WeeklyProgressChart.tsx!} & Biểu
đồ cột 7 ngày \\
\passthrough{\lstinline!src/components/WorkoutCard.tsx!} & Thẻ buổi tập
trình diễn; nội dung hiện đang hard-code \\
\end{longtable}
}

\begin{center}\rule{0.5\linewidth}{0.5pt}\end{center}

\subsection{Badges (Huy hiệu thành
tích)}\label{badges-huy-hiux1ec7u-thuxe0nh-tuxedch}

\subsubsection{Huy hiệu là gì?}\label{huy-hiux1ec7u-luxe0-guxec}

Giống huy chương trong game: khi user đạt mốc nhất định, app \textbf{tự
cấp huy hiệu} và hiển thị sáng trên Dashboard.

\subsubsection{Danh sách 8 huy hiệu mặc
định}\label{danh-suxe1ch-8-huy-hiux1ec7u-mux1eb7c-ux111ux1ecbnh}

{\def\LTcaptype{none} % do not increment counter
\begin{longtable}[]{@{}ll@{}}
\toprule\noalign{}
Tên & Điều kiện đạt \\
\midrule\noalign{}
\endhead
\bottomrule\noalign{}
\endlastfoot
\textbf{Khởi đầu} & Hoàn thành onboarding lần đầu \\
\textbf{Buổi tập đầu} & Hoàn thành 1 buổi tập \\
\textbf{Kiên trì 3 ngày} & Tập liên tiếp 3 ngày (streak \textgreater=
3) \\
\textbf{Tuần vàng} & Tập liên tiếp 7 ngày \\
\textbf{Thép không gỉ} & Tập liên tiếp 30 ngày \\
\textbf{Hydration Pro} & Đạt mục tiêu uống nước trong 1 ngày \\
\textbf{Dinh dưỡng chuẩn} & Đạt mục tiêu calo ăn trong 1 ngày \\
\textbf{Đốt cháy} & Đốt \textgreater= 500 kcal trong một ngày \\
\end{longtable}
}

\subsubsection{Cơ chế cấp huy hiệu (từng
bước)}\label{cux1a1-chux1ebf-cux1ea5p-huy-hiux1ec7u-tux1eebng-bux1b0ux1edbc}

\begin{lstlisting}
Moi lan mo Dashboard
    |
    v
syncUserBadges(userId)
    |
    |-> Thu thap thong ke user:
    |     - Da onboarding chua?
    |     - Tong so buoi tap?
    |     - Streak hien tai?
    |     - Bao nhieu ngay dat muc tieu nuoc/calo?
    |     - Ngay dot calo nhieu nhat?
    |
    |-> So sanh voi tung huy hieu (ham isBadgeEarned)
    |
    |-> Huy hieu du dieu kien ma chua co -> INSERT vao user_badges
    |
    \-> Tra danh sach badge (da dat / chua dat) cho UI
\end{lstlisting}

\textbf{Lưu ý:} Unique constraint
\passthrough{\lstinline!(user\_id, badge\_id)!} bảo đảm mỗi huy hiệu chỉ
được lưu một lần. Service bỏ qua riêng lỗi
\passthrough{\lstinline!23505!} khi hai lần đồng bộ cùng insert một
badge; các lỗi insert khác được ghi cảnh báo.

\subsubsection{Cấu trúc database}\label{cux1ea5u-truxfac-database}

\textbf{Bảng \passthrough{\lstinline!badges!}} --- danh mục huy hiệu
(dùng chung cho mọi user): - \passthrough{\lstinline!code!} --- mã cố
định (vd: \passthrough{\lstinline!streak\_7!}) -
\passthrough{\lstinline!title!}, \passthrough{\lstinline!description!}
--- tên và mô tả hiển thị - \passthrough{\lstinline!criteria\_type!} ---
loại điều kiện (streak, workout\_sessions, \ldots) -
\passthrough{\lstinline!criteria\_value!} --- ngưỡng (vd: 7 ngày)

\textbf{Bảng \passthrough{\lstinline!user\_badges!}} --- huy hiệu user
đã đạt: - \passthrough{\lstinline!user\_id!} +
\passthrough{\lstinline!badge\_id!} --- mỗi cặp chỉ có 1 lần -
\passthrough{\lstinline!earned\_at!} --- thời điểm đạt

\subsubsection{File code liên quan}\label{file-code-liuxean-quan-1}

{\def\LTcaptype{none} % do not increment counter
\begin{longtable}[]{@{}
  >{\raggedright\arraybackslash}p{(\linewidth - 2\tabcolsep) * \real{0.4000}}
  >{\raggedright\arraybackslash}p{(\linewidth - 2\tabcolsep) * \real{0.6000}}@{}}
\toprule\noalign{}
\begin{minipage}[b]{\linewidth}\raggedright
File
\end{minipage} & \begin{minipage}[b]{\linewidth}\raggedright
Vai trò
\end{minipage} \\
\midrule\noalign{}
\endhead
\bottomrule\noalign{}
\endlastfoot
\passthrough{\lstinline!src/services/badgeService.ts!} & Đồng bộ và cấp
huy hiệu \\
\passthrough{\lstinline!src/lib/badgeCalculations.ts!} & Logic kiểm tra
đủ điều kiện \\
\passthrough{\lstinline!src/lib/badgeCatalog.ts!} & Danh mục fallback 4
badge khi không tải được bảng \passthrough{\lstinline!badges!} \\
\passthrough{\lstinline!src/components/BadgeGrid.tsx!} & Hiển thị lưới
huy hiệu \\
\passthrough{\lstinline!src/lib/\_\_tests\_\_/badgeLogic.test.ts!} &
Unit test logic huy hiệu \\
\end{longtable}
}

Danh mục chính trong migration có \textbf{8 badge}, nhưng
\passthrough{\lstinline!FALLBACK\_BADGES!} hiện chỉ có \textbf{4 badge
đầu}. Fallback này cũng vẫn cần các truy vấn Supabase để tính thống kê,
nên không phải chế độ offline hoàn chỉnh.

\begin{center}\rule{0.5\linewidth}{0.5pt}\end{center}

\subsection{Notification (Thông báo \& nhắc
nhở)}\label{notification-thuxf4ng-buxe1o-nhux1eafc-nhux1edf}

\subsubsection{Ba loại thông báo}\label{ba-loux1ea1i-thuxf4ng-buxe1o}

{\def\LTcaptype{none} % do not increment counter
\begin{longtable}[]{@{}
  >{\raggedright\arraybackslash}p{(\linewidth - 4\tabcolsep) * \real{0.2400}}
  >{\raggedright\arraybackslash}p{(\linewidth - 4\tabcolsep) * \real{0.5200}}
  >{\raggedright\arraybackslash}p{(\linewidth - 4\tabcolsep) * \real{0.2400}}@{}}
\toprule\noalign{}
\begin{minipage}[b]{\linewidth}\raggedright
Loại
\end{minipage} & \begin{minipage}[b]{\linewidth}\raggedright
Khi nào bắn
\end{minipage} & \begin{minipage}[b]{\linewidth}\raggedright
Lịch
\end{minipage} \\
\midrule\noalign{}
\endhead
\bottomrule\noalign{}
\endlastfoot
\textbf{Nhắc uống nước} & User bật trong cài đặt & 8 mốc/ngày: 7h, 9h,
11h, 13h, 15h, 17h, 19h, 21h \\
\textbf{Nhắc tập luyện} & User bật trong cài đặt & Đúng giờ
\passthrough{\lstinline!wakeup\_time!} user chọn (vd 6:30 sáng) \\
\textbf{Huy hiệu mới} & Dashboard phát hiện badge đạt mà state trước đó
chưa có & Local notification ngay lập tức \\
\end{longtable}
}

\subsubsection{User bật/tắt ở
đâu?}\label{user-bux1eadttux1eaft-ux1edf-ux111uxe2u}

User có thể mở \passthrough{\lstinline!NotificationSettingsModal!} từ
menu \textbf{3 gạch \textgreater{} Thông báo} hoặc nút chuông trên
Dashboard:

\begin{itemize}
\tightlist
\item
  Công tắc nhắc uống nước
\item
  Công tắc nhắc tập luyện (+ chọn giờ)
\item
  Công tắc thông báo huy hiệu
\end{itemize}

Ba cờ bật/tắt và \passthrough{\lstinline!wakeup\_time!} được lưu trong
\textbf{\passthrough{\lstinline!profiles!}}. Cờ có thể đồng bộ giữa
thiết bị, nhưng quyền OS và id lịch thông báo nằm cục bộ trên từng máy
trong AsyncStorage.

\subsubsection{Luồng bật nhắc tập (ví
dụ)}\label{luux1ed3ng-bux1eadt-nhux1eafc-tux1eadp-vuxed-dux1ee5}

\begin{lstlisting}
User bat "Nhac tap luyen"
    |
    v
App xin quyen thong bao tu he dieu hanh (iOS/Android)
    |
    +-- User tu choi -> bao loi, khong bat
    |
    \-- User dong y
          |
          v
    Lap lich thong bao lap hang ngay dung gio wakeup_time
          |
          v
    Luu id thong bao vao AsyncStorage (de huy sau nay)
          |
          v
    Ghi workout_reminder_enabled = true vao profiles
\end{lstlisting}

Mỗi ngày đúng giờ, điện thoại hiển thị: \textbf{``Giờ tập luyện --- Đã
đến giờ tập!''}

\subsubsection{Lưu ý quan trọng về Expo
Go}\label{lux1b0u-uxfd-quan-trux1ecdng-vux1ec1-expo-go}

\textbf{Thông báo KHÔNG hoạt động trên Expo Go} (app test qua QR code).
Chỉ chạy trên: - \textbf{Development build} (build riêng qua EAS) -
\textbf{Bản release} (APK/IPA cài trên máy)

Code xử lý an toàn: trên Expo Go, module thông báo \textbf{không load}
--- app không crash, chỉ hiển thị thông báo hướng dẫn khi user cố bật
nhắc.

\subsubsection{Khôi phục nhắc khi mở lại
app}\label{khuxf4i-phux1ee5c-nhux1eafc-khi-mux1edf-lux1ea1i-app}

Sau khi đăng nhập và hoàn thành onboarding,
\passthrough{\lstinline!App.tsx!} gọi
\passthrough{\lstinline!syncAllRemindersOnLaunch()!}: - Đọc cờ từ
\passthrough{\lstinline!profiles!} (user đã bật nhắc nước/tập chưa) -
Nếu đã bật -\textgreater{} đặt lại lịch thông báo trên thiết bị

Sau khi cài lại hoặc xóa dữ liệu, app sẽ \textbf{thử} tạo lại lịch từ cờ
server. Việc khôi phục còn phụ thuộc quyền thông báo trên thiết bị và
kết nối Supabase.

\textbf{Giới hạn hiện tại của thông báo badge:}
\passthrough{\lstinline!notifyBadgeEarned()!} không tự xin quyền OS.
Ngoài ra, state badge ban đầu là mảng rỗng nên lần mount Dashboard đầu
tiên có thể coi mọi badge đã đạt là ``mới'' và bắn lại thông báo. Muốn
đúng nghĩa ``chỉ một lần'', cần so sánh với tập badge trước khi gọi
\passthrough{\lstinline!syncUserBadges!} hoặc dùng danh sách id vừa
insert.

\subsubsection{File code liên quan}\label{file-code-liuxean-quan-2}

{\def\LTcaptype{none} % do not increment counter
\begin{longtable}[]{@{}
  >{\raggedright\arraybackslash}p{(\linewidth - 2\tabcolsep) * \real{0.4000}}
  >{\raggedright\arraybackslash}p{(\linewidth - 2\tabcolsep) * \real{0.6000}}@{}}
\toprule\noalign{}
\begin{minipage}[b]{\linewidth}\raggedright
File
\end{minipage} & \begin{minipage}[b]{\linewidth}\raggedright
Vai trò
\end{minipage} \\
\midrule\noalign{}
\endhead
\bottomrule\noalign{}
\endlastfoot
\passthrough{\lstinline!src/services/notificationService.ts!} & Nhắc
tập, prefs, thông báo badge \\
\passthrough{\lstinline!src/services/waterReminderService.ts!} & Nhắc
uống nước (8 mốc) \\
\passthrough{\lstinline!src/components/NotificationSettingsModal.tsx!} &
Giao diện bật/tắt \\
\passthrough{\lstinline!src/components/AppMenuDrawer.tsx!} & Menu mở
modal nhắc nhở \\
\end{longtable}
}

\begin{center}\rule{0.5\linewidth}{0.5pt}\end{center}

\subsection{Luồng hoạt động tổng
hợp}\label{luux1ed3ng-houx1ea1t-ux111ux1ed9ng-tux1ed5ng-hux1ee3p}

\subsubsection{Kịch bản A: User mở app buổi
sáng}\label{kux1ecbch-bux1ea3n-a-user-mux1edf-app-buux1ed5i-suxe1ng}

\begin{lstlisting}
1. Dang nhap -> da onboarding
2. App dong bo nhac nho (neu da bat truoc do)
3. Vao tab Home (Dashboard)
4. Tai song song:
   • So lieu Dashboard (RPC)
   • Dong bo huy hieu (cap moi neu du dieu kien)
   • Doc cai dat thong bao
5. Hien thi: calo, bieu do 7 ngay, huy hieu
6. Neu state truoc chua co badge dat -> thu phat local notification "Ban vua mo khoa ..."
\end{lstlisting}

\subsubsection{Kịch bản B: User hoàn thành buổi
tập}\label{kux1ecbch-bux1ea3n-b-user-houxe0n-thuxe0nh-buux1ed5i-tux1eadp}

\begin{lstlisting}
1. Lam xong bai tap cuoi trong WorkoutDetailScreen
2. Server ghi workout_session + cap nhat streak
3. App tang refreshKey -> Dashboard tai lai
4. Bieu do them cot hom nay, calo tang
5. syncUserBadges co the cap "Buoi tap dau", "Kien tri 3 ngay"...
6. Thu phat thong bao huy hieu (neu bat; hien co rui ro ban lai khi mount)
\end{lstlisting}

\subsubsection{Kịch bản C: User bật nhắc uống
nước}\label{kux1ecbch-bux1ea3n-c-user-bux1eadt-nhux1eafc-uux1ed1ng-nux1b0ux1edbc}

\begin{lstlisting}
1. Menu 3 gach -> Thong bao -> Bat "Nhac uong nuoc"
2. Xin quyen OS -> lap 8 thong bao lap ngay
3. Luu co water_reminder_enabled = true
4. Tu 7h-21h, moi 2 tieng nhan: "Da den gio uong nuoc!"
\end{lstlisting}

\begin{center}\rule{0.5\linewidth}{0.5pt}\end{center}

\subsection{Backend (Database)}\label{backend-database}

\subsubsection{Migration cần chạy}\label{migration-cux1ea7n-chux1ea1y}

File:
\passthrough{\lstinline!supabase/migrations/20260627200000\_dashboard\_badges\_notifications.sql!}

\textbf{Phải chạy SAU} migration workout (tạo bảng
\passthrough{\lstinline!profiles!},
\passthrough{\lstinline!workout\_sessions!},
\passthrough{\lstinline!user\_streaks!}).

Thứ tự đầy đủ:

{\def\LTcaptype{none} % do not increment counter
\begin{longtable}[]{@{}
  >{\raggedright\arraybackslash}p{(\linewidth - 2\tabcolsep) * \real{0.2727}}
  >{\raggedright\arraybackslash}p{(\linewidth - 2\tabcolsep) * \real{0.7273}}@{}}
\toprule\noalign{}
\begin{minipage}[b]{\linewidth}\raggedright
Bước
\end{minipage} & \begin{minipage}[b]{\linewidth}\raggedright
File migration
\end{minipage} \\
\midrule\noalign{}
\endhead
\bottomrule\noalign{}
\endlastfoot
1 &
\passthrough{\lstinline!20260627000000\_nutrition\_water\_health.sql!} \\
2 &
\passthrough{\lstinline!20260627173357\_workout\_onboarding\_atomic\_updates.sql!} \\
3 &
\textbf{\passthrough{\lstinline!20260627200000\_dashboard\_badges\_notifications.sql!}}
\textless- module này \\
4 &
\passthrough{\lstinline!20260701100000\_add\_new\_programs\_v2.sql!} \\
\end{longtable}
}

\subsubsection{Bảng và cột module tạo
thêm}\label{bux1ea3ng-vuxe0-cux1ed9t-module-tux1ea1o-thuxeam}

\textbf{Mở rộng \passthrough{\lstinline!profiles!}:} -
\passthrough{\lstinline!workout\_reminder\_enabled!} --- bật nhắc tập
(mặc định: tắt) -
\passthrough{\lstinline!badge\_notifications\_enabled!} --- thông báo
huy hiệu (mặc định: bật)

\textbf{Bảng mới:} - \passthrough{\lstinline!badges!} --- 8 huy hiệu mặc
định - \passthrough{\lstinline!user\_badges!} --- huy hiệu user đã đạt

\textbf{Hàm RPC mới:} -
\passthrough{\lstinline!get\_dashboard\_summary()!} --- trả JSON tổng
hợp Dashboard

\subsubsection{Bảo mật}\label{bux1ea3o-mux1eadt}

\begin{itemize}
\tightlist
\item
  \textbf{RLS:} User chỉ xem/cấp huy hiệu của chính mình; danh mục
  \passthrough{\lstinline!badges!} chỉ role
  \passthrough{\lstinline!authenticated!} được đọc.
\item
  \textbf{RPC:} Chạy với quyền người gọi
  (\passthrough{\lstinline!SECURITY INVOKER!}), từ chối nếu chưa đăng
  nhập.
\item
  Mỗi user \textbf{không thể} xem streak/calo/huy hiệu của user khác.
\item
  Migration đã có \passthrough{\lstinline!GRANT!} tường minh cho Data
  API, phù hợp với thay đổi mặc định của Supabase năm 2026.
\end{itemize}

\begin{center}\rule{0.5\linewidth}{0.5pt}\end{center}

\subsection{Kiểm thử}\label{kiux1ec3m-thux1eed}

\subsubsection{Lệnh chạy test tự
động}\label{lux1ec7nh-chux1ea1y-test-tux1ef1-ux111ux1ed9ng}

\begin{lstlisting}[language=bash]
npm run test:badges    # test logic huy hieu
npm run test:weekly    # test normalizeWeeklyWorkouts
npm test               # chay tat ca test
npm run typecheck      # kiem tra TypeScript

\end{lstlisting}

\subsubsection{Checklist test thủ
công}\label{checklist-test-thux1ee7-cuxf4ng}

\textbf{Dashboard:} - {[} {]} User mới (0 buổi tập): biểu đồ hiển thị
``Chưa có buổi tập'', không cột mờ - {[} {]} Có buổi tập: cột đúng ngày,
hôm nay highlight xanh - {[} {]} Kéo xuống refresh -\textgreater{} số
liệu cập nhật - {[} {]} Sau khi tập xong -\textgreater{} Dashboard tự
refresh

\textbf{Huy hiệu:} - {[} {]} Tập buổi đầu -\textgreater{} nhận ``Buổi
tập đầu'' - {[} {]} Streak 3 ngày -\textgreater{} nhận ``Kiên trì 3
ngày'' - {[} {]} Huy hiệu đã đạt sáng, chưa đạt mờ

\textbf{Thông báo:} - {[} {]} Expo Go: bật nhắc -\textgreater{} báo lỗi
hướng dẫn, không crash - {[} {]} Dev build: bật nhắc nước
-\textgreater{} nhận đúng 8 mốc - {[} {]} Dev build: bật nhắc tập
-\textgreater{} nhận đúng giờ đã chọn - {[} {]} Đạt huy hiệu mới + bật
thông báo badge -\textgreater{} có thông báo đẩy - {[} {]} Mount lại
Dashboard không bắn lại các badge đã đạt từ trước - {[} {]} Từ chối
quyền OS -\textgreater{} cờ server và lịch local không rơi vào trạng
thái lệch nhau - {[} {]} Kiểm tra dữ liệu 7 ngày quanh 00:00 khi
timezone database khác timezone thiết bị

\textbf{Phạm vi test tự động hiện tại:}
\passthrough{\lstinline!badgeLogic.test.ts!} mới kiểm tra 2 trường hợp
của badge streak (đạt/chưa đạt mốc 3). Chưa có test tự động cho 5 loại
điều kiện còn lại, RPC/fallback, RLS, lịch thông báo và tình huống đồng
bộ cạnh tranh.

\begin{center}\rule{0.5\linewidth}{0.5pt}\end{center}

\subsection{Tóm tắt nhanh}\label{tuxf3m-tux1eaft-nhanh}

{\def\LTcaptype{none} % do not increment counter
\begin{longtable}[]{@{}
  >{\raggedright\arraybackslash}p{(\linewidth - 2\tabcolsep) * \real{0.3913}}
  >{\raggedright\arraybackslash}p{(\linewidth - 2\tabcolsep) * \real{0.6087}}@{}}
\toprule\noalign{}
\begin{minipage}[b]{\linewidth}\raggedright
Câu hỏi
\end{minipage} & \begin{minipage}[b]{\linewidth}\raggedright
Trả lời ngắn
\end{minipage} \\
\midrule\noalign{}
\endhead
\bottomrule\noalign{}
\endlastfoot
Dashboard lấy data từ đâu? & Số liệu chính từ RPC
\passthrough{\lstinline!get\_dashboard\_summary()!}; badge và prefs là
các request riêng \\
Huy hiệu cấp khi nào? & Mỗi lần mở Dashboard,
\passthrough{\lstinline!syncUserBadges!} kiểm tra và cấp tự động \\
Nhắc nhở lưu ở đâu? & Cờ/giờ trong \passthrough{\lstinline!profiles!};
id lịch và quyền OS ở từng thiết bị \\
Tại sao Expo Go không có thông báo? & Giới hạn của Expo Go --- cần
development build \\
Biểu đồ 7 ngày trống? & Đã fix: empty state + luôn đủ 7 ngày từ
server \\
File migration module? &
\passthrough{\lstinline!20260627200000\_dashboard\_badges\_notifications.sql!} \\
\end{longtable}
}

\begin{center}\rule{0.5\linewidth}{0.5pt}\end{center}

\subsection{Giải thích code \& hàm chi
tiết}\label{giux1ea3i-thuxedch-code-huxe0m-chi-tiux1ebft}

Phần này giải thích \textbf{từng hàm quan trọng}: đầu vào, đầu ra, từng
bước xử lý và lý do thiết kế.

\begin{center}\rule{0.5\linewidth}{0.5pt}\end{center}

\subsubsection{\texorpdfstring{\texttt{getDashboardSummary(userId)} ---
\texttt{dashboardService.ts}}{getDashboardSummary(userId) --- dashboardService.ts}}\label{getdashboardsummaryuserid-dashboardservice.ts}

\textbf{Mục đích:} Lấy toàn bộ số liệu Dashboard trong một lần gọi.

RPC tự lấy user từ JWT bằng \passthrough{\lstinline!auth.uid()!}; tham
số \passthrough{\lstinline!userId!} chỉ được dùng khi chuyển sang
fallback client.

\textbf{Code:}

\begin{lstlisting}
export async function getDashboardSummary(userId: string): Promise<DashboardSummary> {
  const { data, error } = await supabase.rpc('get_dashboard_summary');
  if (!error && data) {
    const summary = data as DashboardSummary;
    return {
      ...summary,
      weekly_workouts: normalizeWeeklyWorkouts(summary.weekly_workouts),
    };
  }
  return buildDashboardSummaryClient(userId);
}

\end{lstlisting}

\textbf{Cách hoạt động (từng bước):}

\begin{enumerate}
\def\labelenumi{\arabic{enumi}.}
\tightlist
\item
  \passthrough{\lstinline!supabase.rpc('get\_dashboard\_summary')!} gọi
  hàm PostgreSQL; Supabase client gửi JWT của session hiện tại.
\item
  Khi có \passthrough{\lstinline!data!} và không có
  \passthrough{\lstinline!error!}, app dùng kết quả RPC.
\item
  \passthrough{\lstinline!normalizeWeeklyWorkouts(...)!} chuẩn hóa mảng
  7 ngày trước khi render.
\item
  Khi RPC lỗi hoặc không có data,
  \passthrough{\lstinline!buildDashboardSummaryClient(userId)!} thử đọc
  từng bảng và gom số liệu.
\end{enumerate}

\textbf{Tại sao có fallback?} User vẫn có thể thấy Dashboard khi
migration RPC chưa chạy. Nếu nguyên nhân là mất mạng hoặc lỗi quyền truy
cập chung, fallback cũng có thể thất bại.

\begin{center}\rule{0.5\linewidth}{0.5pt}\end{center}

\subsubsection{\texorpdfstring{\texttt{normalizeWeeklyWorkouts(raw)} ---
\texttt{dashboardService.ts}}{normalizeWeeklyWorkouts(raw) --- dashboardService.ts}}\label{normalizeweeklyworkoutsraw-dashboardservice.ts}

\textbf{Mục đích:} Đảm bảo biểu đồ \textbf{luôn nhận đúng 7 ngày}, kể cả
khi server trả thiếu hoặc date kèm timestamp.

\textbf{Code:}

\begin{lstlisting}
export function normalizeWeeklyWorkouts(raw: unknown): WeeklyWorkoutDay[] {
  const base = buildEmptyWeekly();  // tao san 7 ngay, moi ngay workouts=0, calories=0
  if (!Array.isArray(raw) || raw.length === 0) return base;

  const map = new Map(base.map((day) => [day.date, { ...day }]));
  for (const item of raw) {
    if (!item || typeof item !== 'object') continue;
    const row = item as Record<string, unknown>;
    const date = normalizeDateString(String(row.date ?? '')); // "2026-06-30 00:00:00" -> "2026-06-30"
    if (!map.has(date)) continue;
    map.set(date, {
      date,
      workouts: Number(row.workouts ?? 0),
      calories_burned: Number(row.calories_burned ?? 0),
    });
  }
  return base.map((day) => map.get(day.date) ?? day);
}

\end{lstlisting}

\textbf{Ví dụ minh họa:}

\begin{lstlisting}
Input tu server (thieu ngay, date lech format):
  [{ date: "2026-06-30 00:00:00", workouts: 1, calories_burned: 320 }]

buildEmptyWeekly() tao base 7 ngay: 24/6 -> 30/6, tat ca = 0

Sau khi merge vao map:
  30/6 -> workouts: 1, calories: 320
  cac ngay khac -> giu 0

Output: mang 7 phan tu, dung thu tu ngay
\end{lstlisting}

\textbf{Hàm phụ trợ \passthrough{\lstinline!buildEmptyWeekly()!}:} Dùng
\passthrough{\lstinline!localDateDaysAgo(6 - index)!} --- tính ngày theo
\textbf{timezone máy}, không dùng UTC (tránh lệch ngày khi user ở VN).

\begin{center}\rule{0.5\linewidth}{0.5pt}\end{center}

\subsubsection{\texorpdfstring{\texttt{buildDashboardSummaryClient(userId)}
--- fallback
client}{buildDashboardSummaryClient(userId) --- fallback client}}\label{builddashboardsummaryclientuserid-fallback-client}

\textbf{Mục đích:} Tự gom số liệu khi RPC không dùng được.

\textbf{Cách hoạt động:}

\begin{lstlisting}
Promise.all (chay song song 3 query):
  +-- getProfile(userId)           -> display_name, wakeup_time, co nhac nho
  +-- user_streaks                 -> current_streak, longest_streak
  \-- workout_sessions (7 ngay)    -> workout_date, total_calories

Duyet tung session -> cong don vao weeklyMap theo ngay
Loc session hom nay -> tinh calories_burned_today, workouts_today
Tra object DashboardSummary giong het RPC
\end{lstlisting}

\textbf{Ưu điểm \passthrough{\lstinline!Promise.all!}:} 3 query chạy
\textbf{đồng thời} thay vì tuần tự; tổng thời gian thường gần với query
chậm nhất, nhưng không có bảo đảm cố định là nhanh hơn 3 lần.

\begin{center}\rule{0.5\linewidth}{0.5pt}\end{center}

\subsubsection{\texorpdfstring{\texttt{loadDashboard()} ---
\texttt{DashboardScreen.tsx}}{loadDashboard() --- DashboardScreen.tsx}}\label{loaddashboard-dashboardscreen.tsx}

\textbf{Mục đích:} Hàm trung tâm tải và hiển thị toàn bộ Dashboard.

\textbf{Code:}

\begin{lstlisting}
const loadDashboard = useCallback(async () => {
  const [dashboard, badgeList, prefs] = await Promise.all([
    getDashboardSummary(userId),
    syncUserBadges(userId),
    getNotificationPreferences(userId).catch(() => null),
  ]);

  setSummary(dashboard);
  setBadges((previous) => {
    if (prefs?.badge_notifications_enabled) {
      const previousEarned = new Set(previous.filter(b => b.earned).map(b => b.id));
      for (const badge of badgeList) {
        if (badge.earned && !previousEarned.has(badge.id)) {
          void notifyBadgeEarned('Huy hieu moi', `Ban vua mo khoa "${badge.title}"`);
        }
      }
    }
    return badgeList;
  });
}, [userId]);

\end{lstlisting}

\textbf{Giải thích từng phần:}

\begin{enumerate}
\def\labelenumi{\arabic{enumi}.}
\tightlist
\item
  \textbf{\passthrough{\lstinline!Promise.all!}} --- Tải dashboard,
  badge, prefs \textbf{cùng lúc} (3 request song song).
\item
  \textbf{\passthrough{\lstinline!setSummary(dashboard)!}} - Cập nhật
  greeting, lưới chỉ số nhanh và biểu đồ.
  \passthrough{\lstinline!StatsGrid!} tự tải dữ liệu sức khỏe riêng;
  \passthrough{\lstinline!WorkoutCard!} hiện là nội dung tĩnh.
\item
  \textbf{\passthrough{\lstinline!setBadges!} với callback} --- So sánh
  badge \textbf{cũ} vs \textbf{mới}:

  \begin{itemize}
  \tightlist
  \item
    Tạo \passthrough{\lstinline!Set!} các id badge đã đạt trước đó
    (\passthrough{\lstinline!previousEarned!}).
  \item
    Duyệt \passthrough{\lstinline!badgeList!} mới: nếu badge đạt
    (\passthrough{\lstinline!earned!}) mà \textbf{chưa có trong Set}
    -\textgreater{} gọi \passthrough{\lstinline!notifyBadgeEarned!}.
  \item
    Chỉ bắn thông báo khi
    \passthrough{\lstinline!badge\_notifications\_enabled === true!}.
  \end{itemize}
\item
  \textbf{\passthrough{\lstinline!useCallback(..., [userId])!}} --- Hàm
  chỉ tạo lại khi \passthrough{\lstinline!userId!} đổi, tránh re-render
  vô ích.
\item
  \textbf{\passthrough{\lstinline!refreshKey!}} (prop từ
  \passthrough{\lstinline!App.tsx!}) --- Tăng sau khi hoàn thành buổi
  tập -\textgreater{} \passthrough{\lstinline!useEffect!} gọi lại
  \passthrough{\lstinline!loadDashboard!}.
\end{enumerate}

\textbf{Sai lệch cần lưu ý:} ở lần load đầu,
\passthrough{\lstinline!previous!} là \passthrough{\lstinline![]!}, nên
mọi badge đã đạt đều vắng trong
\passthrough{\lstinline!previousEarned!}. Vì thế code có thể thông báo
lại badge cũ; phép so sánh hiện tại chưa chứng minh badge vừa được
insert trong lần sync này.

\begin{center}\rule{0.5\linewidth}{0.5pt}\end{center}

\subsubsection{\texorpdfstring{\texttt{WeeklyProgressChart} --- biểu đồ
7
ngày}{WeeklyProgressChart --- biểu đồ 7 ngày}}\label{weeklyprogresschart-biux1ec3u-ux111ux1ed3-7-nguxe0y}

\textbf{Input:} \passthrough{\lstinline!days: WeeklyWorkoutDay[]!} từ
\passthrough{\lstinline!summary.weekly\_workouts!}.

\textbf{Logic chính:}

\begin{lstlisting}
const chartDays = useMemo(
  () => (days.length === 7 ? days : buildFallbackDays()),
  [days],
);
const totalWorkouts = chartDays.reduce((sum, day) => sum + day.workouts, 0);
const hasActivity = totalWorkouts > 0;
const maxWorkouts = Math.max(...chartDays.map(d => d.workouts), 1);

\end{lstlisting}

{\def\LTcaptype{none} % do not increment counter
\begin{longtable}[]{@{}
  >{\raggedright\arraybackslash}p{(\linewidth - 2\tabcolsep) * \real{0.4000}}
  >{\raggedright\arraybackslash}p{(\linewidth - 2\tabcolsep) * \real{0.6000}}@{}}
\toprule\noalign{}
\begin{minipage}[b]{\linewidth}\raggedright
Biến
\end{minipage} & \begin{minipage}[b]{\linewidth}\raggedright
Ý nghĩa
\end{minipage} \\
\midrule\noalign{}
\endhead
\bottomrule\noalign{}
\endlastfoot
\passthrough{\lstinline!chartDays!} & Luôn 7 ngày (fallback nếu input
sai) \\
\passthrough{\lstinline!hasActivity!} & \passthrough{\lstinline!false!}
-\textgreater{} hiển thị empty state ``Chưa có buổi tập'' \\
\passthrough{\lstinline!maxWorkouts!} & Ngày tập nhiều nhất = 100\%
chiều cao cột; tối thiểu 1 để tránh chia cho 0 \\
\end{longtable}
}

\textbf{Tính chiều cao cột:}

\begin{lstlisting}
const fillHeight = active
  ? Math.max(12, (day.workouts / maxWorkouts) * CHART_HEIGHT)
  : 0;

\end{lstlisting}

\begin{itemize}
\tightlist
\item
  Ngày có tập (\passthrough{\lstinline!active!}): cột cao tỉ lệ
  \passthrough{\lstinline!workouts / maxWorkouts!}, tối thiểu 12px.
\item
  Ngày không tập: hiển thị chấm nhỏ (\passthrough{\lstinline!barEmpty!})
  thay vì cột mờ 8px (bug cũ đã fix).
\item
  \passthrough{\lstinline!isToday!}: viền xanh lá
  (\passthrough{\lstinline!barTrackToday!}) + nhãn ngày đậm.
\end{itemize}

\begin{center}\rule{0.5\linewidth}{0.5pt}\end{center}

\subsubsection{\texorpdfstring{\texttt{isBadgeEarned(badge,\ stats)} ---
\texttt{badgeCalculations.ts}}{isBadgeEarned(badge, stats) --- badgeCalculations.ts}}\label{isbadgeearnedbadge-stats-badgecalculations.ts}

\textbf{Mục đích:} Hàm \textbf{thuần} (không gọi mạng) --- kiểm tra 1
huy hiệu đã đủ điều kiện chưa. Dễ unit test.

\textbf{Code:}

\begin{lstlisting}
export function isBadgeEarned(badge: Badge, stats: BadgeStats): boolean {
  switch (badge.criteria_type) {
    case 'onboarding':          return stats.onboardingComplete;
    case 'workout_sessions':    return stats.workoutSessions >= badge.criteria_value;
    case 'streak':              return stats.currentStreak >= badge.criteria_value;
    case 'water_goal_days':     return stats.waterGoalDays >= badge.criteria_value;
    case 'nutrition_goal_days': return stats.nutritionGoalDays >= badge.criteria_value;
    case 'calories_burned_day': return stats.maxCaloriesBurnedDay >= badge.criteria_value;
    default:                    return false;
  }
}

\end{lstlisting}

\textbf{Ví dụ:} Badge ``Kiên trì 3 ngày'' có
\passthrough{\lstinline!criteria\_type: 'streak'!},
\passthrough{\lstinline!criteria\_value: 3!}. -
\passthrough{\lstinline!stats.currentStreak = 2!} -\textgreater{}
\passthrough{\lstinline!false!} (chưa đạt). -
\passthrough{\lstinline!stats.currentStreak = 3!} -\textgreater{}
\passthrough{\lstinline!true!} (đạt).

\begin{center}\rule{0.5\linewidth}{0.5pt}\end{center}

\subsubsection{\texorpdfstring{\texttt{loadBadgeStats(userId)} ---
\texttt{badgeService.ts}}{loadBadgeStats(userId) --- badgeService.ts}}\label{loadbadgestatsuserid-badgeservice.ts}

\textbf{Mục đích:} Gom \textbf{một lần} tất cả thống kê cần để kiểm tra
huy hiệu.

\textbf{Các chỉ số thu thập:}

{\def\LTcaptype{none} % do not increment counter
\begin{longtable}[]{@{}
  >{\raggedright\arraybackslash}p{(\linewidth - 2\tabcolsep) * \real{0.4211}}
  >{\raggedright\arraybackslash}p{(\linewidth - 2\tabcolsep) * \real{0.5789}}@{}}
\toprule\noalign{}
\begin{minipage}[b]{\linewidth}\raggedright
Chỉ số
\end{minipage} & \begin{minipage}[b]{\linewidth}\raggedright
Cách tính
\end{minipage} \\
\midrule\noalign{}
\endhead
\bottomrule\noalign{}
\endlastfoot
\passthrough{\lstinline!onboardingComplete!} &
\passthrough{\lstinline!profiles.onboarding\_completed === true!} \\
\passthrough{\lstinline!workoutSessions!} &
\passthrough{\lstinline!COUNT(*)!} từ
\passthrough{\lstinline!workout\_sessions!} \\
\passthrough{\lstinline!currentStreak!} &
\passthrough{\lstinline!user\_streaks.current\_streak!} \\
\passthrough{\lstinline!waterGoalDays!} & Đếm ngày
\passthrough{\lstinline!water\_ml >= water\_goal\_ml!} (+ hôm nay nếu
đạt) \\
\passthrough{\lstinline!nutritionGoalDays!} & Đếm ngày
\passthrough{\lstinline!calories\_consumed >= daily\_calorie\_goal!} \\
\passthrough{\lstinline!maxCaloriesBurnedDay!} & Gom calo theo
\passthrough{\lstinline!workout\_date!}, lấy max \\
\end{longtable}
}

\textbf{Lưu ý:} Hôm nay được tính riêng qua
\passthrough{\lstinline!getDailyNutrition!} /
\passthrough{\lstinline!getTodayWater!} vì có thể chưa có dòng trong
bảng daily (chưa ghi trong ngày).

\passthrough{\lstinline!nutritionGoalDays!} so sánh mọi ngày lịch sử với
\textbf{mục tiêu calo hiện tại} của profile, không lưu snapshot mục tiêu
theo từng ngày. Nếu user đổi mục tiêu, số ngày đạt chuẩn trong quá khứ
có thể thay đổi.

\begin{center}\rule{0.5\linewidth}{0.5pt}\end{center}

\subsubsection{\texorpdfstring{\texttt{syncUserBadges(userId)} --- cấp
huy hiệu tự
động}{syncUserBadges(userId) --- cấp huy hiệu tự động}}\label{syncuserbadgesuserid-cux1ea5p-huy-hiux1ec7u-tux1ef1-ux111ux1ed9ng}

\textbf{Code:}

\begin{lstlisting}
// Rut gon tu badgeService.ts; code that co kiem tra badgesError.
export async function syncUserBadges(userId: string): Promise<BadgeWithStatus[]> {
  const [{ data: badges, error: badgesError }, stats, { data: earned }] = await Promise.all([
    supabase.from('badges').select('*').order('sort_order', { ascending: true }),
    loadBadgeStats(userId),
    supabase.from('user_badges').select('badge_id').eq('user_id', userId),
  ]);

  if (badgesError || !badges?.length) return getBadgesWithStatus(userId);

  const earnedIds = new Set((earned ?? []).map(r => r.badge_id));
  const toAward = badges.filter(b => !earnedIds.has(b.id) && isBadgeEarned(b, stats));

  if (toAward.length > 0) {
    const { error } = await supabase.from('user_badges').insert(
      toAward.map(b => ({ user_id: userId, badge_id: b.id })),
    );
    if (error && error.code !== '23505') console.warn(error.message);
  }
  return getBadgesWithStatus(userId);
}

\end{lstlisting}

\textbf{Luồng:}

\begin{lstlisting}
1. Tai song song: danh muc badge + stats user + badge da dat
2. earnedIds = Set cac badge_id da co trong user_badges
3. toAward = badge chua co VA isBadgeEarned() = true
4. INSERT toAward vao user_badges
5. Tra getBadgesWithStatus() — danh sach day du kem earned/earned_at
\end{lstlisting}

\textbf{Tại sao bỏ qua lỗi \passthrough{\lstinline!23505!}?} Hai lần gọi
\passthrough{\lstinline!syncUserBadges!} đồng thời có thể cùng cố insert
một badge. Tuy nhiên đây là insert theo batch; nếu batch gặp conflict,
cần kiểm thử để chắc các badge khác trong batch không bị bỏ lỡ. Cách rõ
ràng hơn là dùng upsert/\passthrough{\lstinline!ON CONFLICT DO NOTHING!}
phù hợp với policy.

\begin{center}\rule{0.5\linewidth}{0.5pt}\end{center}

\subsubsection{\texorpdfstring{\texttt{getNotifications()} --- lazy-load
an toàn Expo
Go}{getNotifications() --- lazy-load an toàn Expo Go}}\label{getnotifications-lazy-load-an-touxe0n-expo-go}

\textbf{Mục đích:} Chỉ load \passthrough{\lstinline!expo-notifications!}
khi \textbf{không} chạy trên Expo Go.

\begin{lstlisting}
const isExpoGo = Constants.executionEnvironment === ExecutionEnvironment.StoreClient;

async function getNotifications(): Promise<NotificationsModule | null> {
  if (isExpoGo) return null;  // Expo Go -> khong dung notification

  if (notificationsModule === undefined) {
    try {
      notificationsModule = await import('expo-notifications');  // import dong, lan dau
      notificationsModule.setNotificationHandler({ /* hien alert + am thanh */ });
    } catch {
      notificationsModule = null;
    }
  }
  return notificationsModule;
}

\end{lstlisting}

\textbf{Pattern lazy-load:} - \passthrough{\lstinline!undefined!} = chưa
thử load. - \passthrough{\lstinline!null!} = đã thử import nhưng lỗi. -
\passthrough{\lstinline!NotificationsModule!} = đã load thành công và
được cache. - Trên Expo Go, hàm trả \passthrough{\lstinline!null!} sớm;
biến cache vẫn giữ nguyên.

-\textgreater{} Mọi hàm notification gọi
\passthrough{\lstinline!getNotifications()!} trước; nếu
\passthrough{\lstinline!null!} -\textgreater{} return sớm, \textbf{không
crash}.

\begin{center}\rule{0.5\linewidth}{0.5pt}\end{center}

\subsubsection{\texorpdfstring{\texttt{updateNotificationPreferences(userId,\ prefs)}
--- bật/tắt
nhắc}{updateNotificationPreferences(userId, prefs) --- bật/tắt nhắc}}\label{updatenotificationpreferencesuserid-prefs-bux1eadttux1eaft-nhux1eafc}

\textbf{Mục đích:} Cập nhật cài đặt nhắc nhở --- vừa lập lịch trên máy,
vừa ghi cờ lên server.

\textbf{Luồng xử lý:}

\begin{lstlisting}
1. Doc prefs hien tai: getNotificationPreferences(userId)
2. Merge: next = { ...current, ...prefs }

3. Neu bat nhac nuoc (true):
     -> enableWaterReminders() — xin quyen + lap 8 thong bao lap ngay
   Neu tat (false):
     -> disableWaterReminders() — huy tat ca id da luu trong AsyncStorage

4. Neu bat nhac tap (true):
     -> requestPermissions() — xin quyen OS
     -> (Android) tao notification channel
     -> scheduleWorkoutReminder(wakeup_time) — lap lich CALENDAR lap ngay
   Neu tat -> cancelWorkoutReminder()
   Neu doi gio khi dang bat -> scheduleWorkoutReminder(gio moi)

5. Ghi 3 co boolean va `wakeup_time` len `profiles` (Supabase)
6. Tra next prefs
\end{lstlisting}

Luồng này không phải transaction giữa OS và database: lịch local được
đổi trước, rồi mới update \passthrough{\lstinline!profiles!}. Nếu update
Supabase thất bại, lịch trên máy và cờ server có thể lệch nhau.

\textbf{\passthrough{\lstinline!scheduleWorkoutReminder(wakeupTime)!}
chi tiết:}

\begin{lstlisting}
await cancelWorkoutReminder();           // huy lich cu (tranh trung)
const { hour, minute } = parseWakeupTime(wakeupTime);  // "06:30" -> {6, 30}
const id = await Notifications.scheduleNotificationAsync({
  content: { title: 'Giờ tập luyện ', body: '...', sound: true },
  trigger: { type: CALENDAR, hour, minute, repeats: true },  // lap moi ngay
});
await AsyncStorage.setItem(WORKOUT_REMINDER_ID_KEY, id);  // luu id de huy sau

\end{lstlisting}

\begin{center}\rule{0.5\linewidth}{0.5pt}\end{center}

\subsubsection{\texorpdfstring{\texttt{syncAllRemindersOnLaunch(userId)}
--- khôi phục khi mở
app}{syncAllRemindersOnLaunch(userId) --- khôi phục khi mở app}}\label{syncallremindersonlaunchuserid-khuxf4i-phux1ee5c-khi-mux1edf-app}

\textbf{Gọi từ:} \passthrough{\lstinline!App.tsx!} sau khi xác nhận user
đã onboarding.

\begin{lstlisting}
export async function syncAllRemindersOnLaunch(userId: string): Promise<void> {
  if (isExpoGo) return;

  await syncWaterRemindersOnLaunch(userId);   // doc co water -> dat lai 8 moc neu bat

  const prefs = await getNotificationPreferences(userId);
  if (prefs.workout_reminder_enabled) {
    const granted = await requestPermissions();
    if (granted) await scheduleWorkoutReminder(prefs.wakeup_time);
  }
}

\end{lstlisting}

\textbf{Tại sao cần?} Khi id trong AsyncStorage không còn, app có thể
dựa vào cờ server để thử đặt lại lịch. Kết quả vẫn phụ thuộc quyền OS và
kết nối mạng; lỗi trong \passthrough{\lstinline!App.tsx!} hiện chỉ được
ghi bằng \passthrough{\lstinline!console.error!}.

\begin{center}\rule{0.5\linewidth}{0.5pt}\end{center}

\subsubsection{\texorpdfstring{RPC \texttt{get\_dashboard\_summary()}
--- SQL trên
server}{RPC get\_dashboard\_summary() --- SQL trên server}}\label{rpc-get_dashboard_summary-sql-truxean-server}

\textbf{Phần quan trọng nhất --- snapshot 7 ngày:}

\begin{lstlisting}[language=SQL]
from generate_series(v_today - 6, v_today, interval '1 day') as day
left join lateral (
  select count(*)::int as count, coalesce(sum(total_calories),0) as calories
  from public.workout_sessions
  where user_id = v_user_id and workout_date = day::date
) ws on true

\end{lstlisting}

\textbf{Giải thích:} - \passthrough{\lstinline!generate\_series!} tạo
\textbf{7 ngày liên tiếp} (hôm nay - 6 -\textgreater{} hôm nay). -
\passthrough{\lstinline!LEFT JOIN LATERAL!} với subquery đếm session
theo từng ngày. - Ngày \textbf{không có buổi tập} vẫn xuất hiện với
\passthrough{\lstinline!workouts = 0!},
\passthrough{\lstinline!calories\_burned = 0!}. - Việc giữ đủ ngày đến
từ \passthrough{\lstinline!generate\_series!} kết hợp subquery aggregate
luôn trả một dòng; không nên khẳng định mọi
\passthrough{\lstinline!INNER JOIN!} đều làm mất ngày nếu chưa xét hình
dạng subquery. - \passthrough{\lstinline!v\_today := current\_date!}
theo timezone database. Cần cấu hình/chuẩn hóa timezone để khớp với ngày
local trên thiết bị.

\textbf{Bảo mật:} -
\passthrough{\lstinline!v\_user\_id := (select auth.uid())!} --- chỉ lấy
data của user đang đăng nhập. -
\passthrough{\lstinline!SECURITY INVOKER!} --- chạy với quyền user, RLS
vẫn áp dụng. -
\passthrough{\lstinline!if v\_user\_id is null then raise exception!}
--- từ chối nếu chưa đăng nhập.

\begin{center}\rule{0.5\linewidth}{0.5pt}\end{center}

\subsubsection{Tóm tắt hàm}\label{tuxf3m-tux1eaft-huxe0m}

\begin{itemize}
\tightlist
\item
  \passthrough{\lstinline!getDashboardSummary!}
  (\passthrough{\lstinline!dashboardService!}): nhận
  \passthrough{\lstinline!userId!}, trả
  \passthrough{\lstinline!DashboardSummary!}; gọi khi mở hoặc refresh
  Dashboard.
\item
  \passthrough{\lstinline!normalizeWeeklyWorkouts!}
  (\passthrough{\lstinline!dashboardService!}): nhận
  \passthrough{\lstinline!raw: unknown!}, trả mảng 7 ngày; chạy sau RPC.
\item
  \passthrough{\lstinline!loadDashboard!}
  (\passthrough{\lstinline!DashboardScreen!}): tải song song summary,
  badge và prefs rồi cập nhật state.
\item
  \passthrough{\lstinline!isBadgeEarned!}
  (\passthrough{\lstinline!badgeCalculations!}): nhận badge và stats,
  trả boolean cho từng điều kiện.
\item
  \passthrough{\lstinline!syncUserBadges!}
  (\passthrough{\lstinline!badgeService!}): kiểm tra/cấp badge khi
  Dashboard load.
\item
  \passthrough{\lstinline!getNotificationPreferences!} và
  \passthrough{\lstinline!updateNotificationPreferences!}
  (\passthrough{\lstinline!notificationService!}): đọc/ghi cài đặt
  modal.
\item
  \passthrough{\lstinline!syncAllRemindersOnLaunch!}
  (\passthrough{\lstinline!notificationService!}): thử khôi phục lịch
  sau login và onboarding.
\item
  \passthrough{\lstinline!notifyBadgeEarned!}
  (\passthrough{\lstinline!notificationService!}): phát local
  notification tức thời; hiện chưa tự xin quyền.
\item
  \passthrough{\lstinline!get\_dashboard\_summary!} (SQL RPC): lấy user
  từ JWT và trả \passthrough{\lstinline!jsonb!}.
\end{itemize}

\begin{center}\rule{0.5\linewidth}{0.5pt}\end{center}

\emph{Tài liệu giải thích module Dashboard, Notification \& Badges ---
FitLife. Phiên bản dễ hiểu + giải thích code cho nghiệm thu và báo cáo
nhóm.}

\end{document}
